\documentclass[10pt,a4paper]{amsart}
\usepackage[margin=19mm]{geometry}
\usepackage{amsmath,amssymb,mathtools}
\usepackage[scr=rsfs,cal=euler]{mathalfa}
\usepackage{xfrac}
\usepackage[unicode,hidelinks,bookmarks=false]{hyperref}
\newcommand{\ie}{i.e.\ }
\newcommand{\cf}{cf.\ }
\newcommand{\resp}{respectively\ }
\newcommand{\Subsection}{\subsection}
\providecommand{\nobreakdash}{\nobreak\hbox{-}}
\setlength{\emergencystretch}{2em}
\multlinegap0pt
\usepackage[shortlabels]{enumitem}
\usepackage{amssymb}
\usepackage{mathtools}
\usepackage{bm}
\usepackage{xcolor}
\usepackage{float}
\usepackage{tcolorbox}
\newtheorem{theorem}{Theorem}
\newcommand{\E}{\mathbb{E}}
\newcommand{\R}{\mathbb{R}}
\newcommand{\dif}{\,\mathrm{d}}
\newcommand{\norm}[1]{\left\lVert #1 \right\rVert}
\title{\textbf{Joint Lyapunov Certificates for $K$-Agent\\ Generative AI Governance}\\[0.5em] \large Stochastic Stability, Emergent Ensemble Risk,\\ and Zero-Knowledge Governance Attestation}
\author{Sriram Nagaraj}
\date{Source: arXiv:2608.09087v1 (CC BY 4.0)}
\begin{document}
\maketitle

\noindent\emph{Source and licence.} \textbf{Joint Lyapunov Certificates for $K$-Agent\\ Generative AI Governance}\\[0.5em] \large Stochastic Stability, Emergent Ensemble Risk,\\ and Zero-Knowledge Governance Attestation. Version arXiv:2608.09087v1 (CC BY 4.0), distributed by arXiv under the Creative Commons Attribution 4.0 International licence, http://creativecommons.org/licenses/by/4.0/, as recorded in the arXiv licence metadata for this exact version and shown on the abstract page https://arxiv.org/abs/2608.09087v1. Full licence notice: \texttt{licenses/arxiv-2608.09087-CC-BY-4.0.txt}.\par\smallskip

\noindent\textit{Editorial scope.} Counted unit: the article's theorem thm:emergent (source0.tex lines 1077--1102). The article's complete original proof of it is delivered with it (source0.tex lines 1104--1158, one \texttt{proof} environment of 55 lines). No mathematics is written, repaired or re-derived: the statement and the proof are the article's own lines, in the article's own notation, and no retained line is substituted or re-flowed.  No further source-local context is needed: every cross-reference of every retained line resolves inside the two delivered spans.  Nothing was omitted from any retained span, so the package carries no omission record.  No retained line contains a citation, so the wrapper carries no bibliography and no bibliography entry.  No retained line contains a figure environment, an image-inclusion command or drawing code, so no image is delivered and the package declares no asset dependency.  Editorial formatting only, in addition: an AMS \texttt{amsart} preamble that restates the article's own declarations and macros used by the retained lines, namely the environments \texttt{theorem} on separate counters, together with the standard packages those lines need.  The retained environments print in this document as Theorem 1 in order of appearance, whereas the article prints the same items with the numbering of its own full document.  The complete native source of the article as distributed in the arXiv e-print for this exact version is bundled unchanged as \texttt{sources/207213/source0.tex}.
\begin{theorem}[Emergent ensemble risk]
  \label{thm:emergent}
  Consider the decoupled system ($\gamma = 0$, $\Phi = 0$) with $K$ agents
  and suppose a hidden drift $h \in \R^d$ is added to every agent's
  dynamics, so the true drift of agent $k$ is
  $\mu^k(W^k) = -\alpha_{\mathrm{self}} W^k + h$.
  Define the honest noise floor
  $V^\infty_{\mathrm{hon}} = Kd\sigma_0^2/(4\alpha_{\mathrm{self}})$.

  If the drift magnitude satisfies
  \begin{equation}
    \frac{K}{2\alpha_{\mathrm{self}}^2}\norm{h}^2
    = \delta \cdot V^\infty_{\mathrm{hon}}
    \quad\text{with}\quad 0 < \delta < \frac{\theta_{\mathrm{ind}}}
    {V^\infty_{\mathrm{hon}} / K} - 1,
    \label{eq:h_condition}
  \end{equation}
  then:
  \begin{enumerate}[label=(\roman*)]
    \item No individual agent is flagged: for every $k$,
      $\E_\pi[V^k] < \theta_{\mathrm{ind}}$.
    \item The joint monitor flags the system:
      $V_{\mathrm{joint}} = \sum_k \E_\pi[V^k] >
      (1 + \delta/2) \cdot V^\infty_{\mathrm{hon}}$.
  \end{enumerate}
\end{theorem}

\begin{proof}
  \textbf{Stationary distribution under drift.}
  With hidden drift $h$, the $i$-th component of agent $k$'s weights
  evolves as
  \[
    \dif W^k_{t,i} = (-\alpha_{\mathrm{self}} W^k_{t,i} + h_i)\,\dif t
    + \sigma_0\,\dif B^k_{t,i}.
  \]
  This is an OU process with mean $h_i/\alpha_{\mathrm{self}}$, with
  stationary distribution
  \[
    \mathcal{N}\!\left(h_i/\alpha_{\mathrm{self}},\; \sigma_0^2/(2\alpha_{\mathrm{self}})\right).
  \]

  \textbf{Stationary per-agent second moment.}
  \[
    \E_\pi[V^k] = \frac{1}{2}\sum_{i=1}^d \E_\pi[(W^k_i)^2]
    = \frac{1}{2}\sum_{i=1}^d \left[\frac{h_i^2}{\alpha_{\mathrm{self}}^2}
    + \frac{\sigma_0^2}{2\alpha_{\mathrm{self}}}\right]
    = \frac{\norm{h}^2}{2\alpha_{\mathrm{self}}^2}
    + \frac{d\sigma_0^2}{4\alpha_{\mathrm{self}}}.
  \]

  \textbf{Part (i): Individual check.}
  The per-agent honest mean is
  $\E_\pi^{\mathrm{hon}}[V^k] = d\sigma_0^2/(4\alpha_{\mathrm{self}})
  = V^\infty_{\mathrm{hon}}/K$.
  The individual threshold is $\theta_{\mathrm{ind}} = 2V^\infty_{\mathrm{hon}}/K$.
  Agent $k$ is flagged only if
  \[
    \E_\pi[V^k] = \frac{\norm{h}^2}{2\alpha_{\mathrm{self}}^2}
    + \frac{V^\infty_{\mathrm{hon}}}{K}
    > \frac{2V^\infty_{\mathrm{hon}}}{K},
  \]
  i.e., if $\norm{h}^2 > 2\alpha_{\mathrm{self}}^2 V^\infty_{\mathrm{hon}} / K$.
  By condition \eqref{eq:h_condition},
  $\norm{h}^2 = 2\delta\alpha_{\mathrm{self}}^2 V^\infty_{\mathrm{hon}} / K$,
  and the right-hand condition requires
  $\delta < (\theta_{\mathrm{ind}}/V^\infty_{\mathrm{hon}})\cdot K - 1
  = 1$.  Under this condition, no agent is individually flagged.

  \textbf{Part (ii): Joint check.}
  The joint stationary value is
  \[
    V_{\mathrm{joint}} = K \cdot \E_\pi[V^k]
    = \frac{K\norm{h}^2}{2\alpha_{\mathrm{self}}^2}
    + V^\infty_{\mathrm{hon}}.
  \]
  With the choice of $h$ in \eqref{eq:h_condition},
  the first term equals $\delta V^\infty_{\mathrm{hon}}$, so
  $V_{\mathrm{joint}} = (1 + \delta) V^\infty_{\mathrm{hon}}$.
  The joint threshold is $(1+\delta/2) V^\infty_{\mathrm{hon}} <
  (1+\delta) V^\infty_{\mathrm{hon}} = V_{\mathrm{joint}}$, so the
  joint monitor fires.
\end{proof}

\end{document}
